%% ============================================================
%% Chapter 21: Unfinishable Games and Temporal Extraction
%% ============================================================

\chapter{Unfinishable Games and Temporal Extraction}
\label{ch:ch21-unfinishable-games}
\section{Two Critiques, One Vocabulary}

Chapter~\ref{ch:ch19-ledgers-without-value} examined how a monetized game inverts the developmental function of difficulty, converting a training environment into an extraction mechanism through engineered scarcity and variable-ratio reinforcement. This chapter examines a second, analytically distinct extraction mechanism present in the same broad genre, one that requires no purchase at all: structural unfinishability, the removal of a participant's unilateral control over when disengagement from a persistent, real-time, alliance-based game is safe.

The last decade of criticism directed at exploitative game design has converged almost entirely on monetization systems that use randomized rewards to convert player spending into a variable-ratio reinforcement schedule structurally similar to a slot machine. This body of criticism is substantial and its harms are real, though the international regulatory picture is considerably more fragmented than a simple narrative of bans would suggest \cite{xiao2021}: Belgium's gaming commission found certain paid loot box implementations to constitute gambling, while the Netherlands' comparable finding was later overturned on administrative appeal, and the United Kingdom's government-commissioned review concluded that an association between loot box spending and gambling-related harm existed but that a causal relationship had not been established. None of this requires the game to run in real time. A turn-based game, or one with no live multiplayer component at all, can implement an identical loot box system with identical monetary extraction and identical gambling-adjacent psychology. The mechanism is orthogonal to the game's temporal structure, and this orthogonality means a second mechanism, tied specifically to real-time synchronous structure, cannot be a mere variant or intensification of the first. It has to be examined on its own terms, with its own formal apparatus, rather than folded into the borrowed vocabulary of addiction, which tends to relocate the explanatory burden onto the player's self-control rather than onto the properties of the system that make stopping costly.

\section{Formalizing Structural Unfinishability}

A precise definition needs to separate four phenomena an informal account of ``real-time games'' tends to run together: persistent simulation, globally distributed participation, consequential events occurring during a participant's absence, and social responsibility for preventing collective losses. No one of these alone produces the mechanism this chapter is concerned with. A persistent world can permit harmless absence. A globally distributed player base can coordinate asynchronously. Consequential events can be bounded by protection mechanics. Responsibility can be spread widely enough that no single participant remains continuously accountable. The mechanism arises specifically from their conjunction.

\begin{definition}[Structural unfinishability]
A game is structurally unfinishable for participant $i$ when consequential state changes can occur during $i$'s absence, when their timing is substantially controlled by other players or the server rather than by $i$, when the resulting losses are difficult to reverse, and when responsibility for preventing those losses is socially attributed to $i$.
\end{definition}

This makes unfinishability relational rather than a fixed property of the software: the same game can be an easily stoppable pastime for an ordinary member and an effectively continuous obligation for an alliance leader, because $R_i(t)$, the responsibility attributed to participant $i$, and $U_i(t)$, the difficulty of substituting another participant for $i$, both vary by role. A hazard or event-intensity term $\lambda_i(t)$, the expected rate of consequential events conditional on $i$'s absence, is estimable in principle from server logs. Let $L_i(t)$ denote the expected loss to $i$ conditional on an event occurring while absent, let $R_i(t)$ denote socially attributed responsibility, and let $U_i(t)$ denote the difficulty of substitution. Define $i$'s temporal exposure over an interval $T$ as
\[
X_i(T) = \frac{1}{|T|} \int_T \lambda_i(t)\, L_i(t)\, R_i(t)\, U_i(t)\, dt.
\]
A safe disengagement interval $I \subseteq T$ exists when $\int_I \lambda_i(t) L_i(t) R_i(t) U_i(t)\, dt \leq \varepsilon$ for some tolerable threshold. This decomposition separates four independent levers, exposure, cost, attributed responsibility, and irreplaceability, and a design intervention can reduce unfinishability by attacking any one of the four without needing to touch the others. Structural unfinishability is not the literal absence of any interval of low exposure; it is the scarcity, or the participant's inability to unilaterally guarantee the existence, of a sufficiently long interval satisfying the inequality. This makes the concept comparative and, in principle, measurable, and it draws a necessary line between persistence and unfinishability: a voxel-building server can operate continuously all night without $L_i(t)$ or $R_i(t)$ rising appreciably for any particular player, because nothing consequential and attributable to that player is at stake. Unfinishability names the specific condition in which the world's continuous operation is coupled to a specific participant's exposure rather than decoupled from it.

\section{The Circadian Vulnerability}

A further asymmetry explains why global synchronicity in particular, rather than mere real-time operation, aggravates $\lambda_i(t)$ for a globally distributed alliance. Human availability is approximately periodic on a roughly twenty-four hour cycle tied to local time zone. Let $a_j(t)$ denote attacker $j$'s availability, and let $A(t) = \sum_j a_j(t)$ denote the aggregate availability of the full adversarial population. When attackers are concentrated in a single time zone, $A(t)$ inherits a strong daily rhythm; as the population's time-zone dispersion increases, the individual peaks and troughs of $a_j(t)$ increasingly fall at different points in the day, and $\operatorname{Var}_t[A(t)]$ falls correspondingly, even while each individual attacker retains their own strongly periodic availability. In the limit of a sufficiently large and well-distributed adversary population, $A(t)$ approaches a roughly constant level across the full cycle, while a defending player retains their own strongly periodic availability, low during local sleeping hours. The defender's least-available hours are, in expectation, no less contested than any other hours. Circadian absence, a biological constant no game design choice can remove, is thereby converted into a strategically exploitable regularity in $\lambda_i(t)$ purely as a function of the adversarial population's geographic dispersion, without requiring any single attacker to act with deliberate intent, yielding a testable prediction: the severity of the effect should scale with the variance reduction in $A(t)$ as dispersion increases.

\section{Delegability}

Delegability, formalized as low $U_i(t)$, determines whether global synchronicity actually forces continuous individual availability or can instead be absorbed by a group. If responsibility for responding to a threat can be handed seamlessly to another participant, no single individual need remain continuously available even though $\lambda_i(t)$ itself never falls to zero. The configuration this chapter is concerned with therefore requires a specific conjunction of high $\lambda_i$, high $L_i$, high $R_i$, and high $U_i$ simultaneously; removing any one, through protected rest periods, reliable delegation, bounded and recoverable losses, or a genuinely redundant responsibility structure, moves a game outside the category even if it remains globally synchronous. This predicts which superficially similar games should escape the mechanism, not only which ones exhibit it.

\section{From Session to Membership}

A server clock alone cannot make a person wake at an unscheduled hour; what supplies the compulsion, formally, is $R_i(t)$, a social quantity assigned by other participants who are waiting, depending on the player's presence, or standing to suffer a consequence alongside them. This is clearest by contrast with an older model of multiplayer participation, real-time strategy titles played over a modem, in which multiplayer was an optional session: players deliberately connected, played a bounded match, and disconnected, with no persistent structure surviving the session's end. A modern alliance-based persistent game replaces the match with membership:
\[
\text{multiplayer as an optional session} \longrightarrow \text{persistent social membership}.
\]
Structural unfinishability, in the fullest sense intended here, is what results when membership of this kind is combined with the exposure structure above: the game has stopped being a session one opts into and become a social position one occupies continuously, with occupancy itself carrying consequences for others.

$R_i(t)$ cannot arise from a server clock alone. A clock can make $\lambda_i(t)$ nonzero at a given hour, but it cannot, by itself, attach responsibility for the consequences of that hour to any specific person. That attribution requires durable relational infrastructure: an alliance that remembers who holds which rank, a chat channel through which expectations are communicated and enforced, a reputation system through which absence becomes visible to others. Once a game supplies persistent identity, standing group membership, user-to-user communication, and reputation, the software is performing social-platform functions in addition to game functions, and this compounds a familiar mechanism, absence made socially legible and costly, with a second, mechanical one ordinary social media lacks: a participant's assets, territory, or alliance standing can be materially changed by other participants specifically because they were absent.

\section{Extraction Requires a Beneficiary}

Burden alone does not establish extraction in any economically meaningful sense; extraction requires showing that the burden's product is captured by someone other than the person bearing it. Unfinishability is the structural condition analyzed above; extraction occurs only when some actor captures value produced under that condition, and the two are logically separable. A nonprofit, open-source persistent game could in principle exhibit exactly the same $\lambda_i$, $L_i$, $R_i$, $U_i$ profile, the same structural unfinishability, without any publisher capturing economic surplus from the resulting coordination labor, because no publisher exists to capture it. Unfinishability would remain a real burden, but calling it extraction in that case would be a category error.

The relevant claim for a commercially operated game is therefore a distinct, additional hypothesis. Let $V_{\mathrm{social}}$ denote the value of the player-produced organization that coordination labor generates, recruitment, scheduling, dispute resolution, and general cohesion of an active player base, and consider the hypothesis $\partial\,\mathrm{Retention}/\partial\,V_{\mathrm{social}} > 0$: that a publisher's retention rises with the value of this uncompensated player-produced organization. Stated this way, extraction becomes an empirically testable economic hypothesis about a specific game and publisher, rather than a term the argument's own vocabulary presupposes:
\[
\text{player coordination labor} \longrightarrow \text{persistent social content}
\]
\[
\longrightarrow \text{retention and monetizable engagement}.
\]
An alliance leader who recruits and retains members, schedules events, moderates disputes, and manufactures the rivalries that make ongoing play compelling is producing exactly the kind of durable social content that keeps a persistent game populated and worth continuing to operate for everyone else in that alliance \cite{levine2008,nardi2010,yee2006}. This labor is not compensated by the publisher. Whether the hypothesis holds for a given commercial title is an empirical matter this chapter does not settle; what it establishes is the more limited claim that appropriation, a publisher benefiting functionally from an active, well-organized population, is not automatically the same thing as extraction, deliberate capture of a specific economic surplus, and that a commercially operated persistent game exhibiting structural unfinishability should be examined for the latter rather than assumed to practice it merely because the former is present.

\section{Case Study: \emph{Last Shelter: Survival}}

\emph{Last Shelter: Survival} provides a concrete, bounded illustration of this architecture, presented here as an example, not as evidence that every game in its genre produces identical effects. The developer's own listing describes the game as a real-time strategy title built around persistent base development, alliance formation, and continuous conflict between players worldwide \cite{lastshelter_appstore,lastshelter_playstore}, and the storefront listing discloses loot boxes, in-app purchases, and persistent messaging, with an age rating of thirteen and older. The game's competitive structure includes a server-versus-server ``State vs State'' competition that a community-maintained guide describes as running on a roughly forty-eight-hour cycle \cite{lastshelter_swvs,lastshelter_appgamer}, during which an alliance's collective standing depends on continuous, coordinated participation across the event's full duration. A partial, individually scoped mitigation exists in the form of a temporary defensive shield subject to a reactivation cooldown, but this does not by itself relieve the coordination requirement borne by alliance leadership during a live, alliance-wide event.

An ordinary alliance member may have comparatively low $R_i(t)$ and low $U_i(t)$: their individual absence is substitutable. An alliance leader, diplomat, or event coordinator occupies a different position on both quantities, with responsibility for war declarations and defensive scheduling commonly and visibly attributed to specific ranked roles, and, particularly in smaller alliances, the difficulty of substituting another coordinator during an active event may be high. Unfinishability, on this reading, is not a claim about the software in the abstract but a claim about a specific, identifiable role within it, exactly the relational property the formal definition was built to capture. Official storefront listings establish $\lambda_i(t)$ and, more weakly, $L_i(t)$ reasonably well; player-produced community guides are reasonable evidence that particular mechanics exist and function as described; neither source establishes $R_i(t)$ or $U_i(t)$ directly, since both are social quantities that would properly require observation of actual alliance communications or survey data to measure. What the case study shows is that \emph{Last Shelter: Survival} instantiates the technical preconditions under which the proposed mechanism can arise, while whether particular players in particular roles actually cross a specified $X_i$ threshold remains an empirical question the case study does not itself answer:
\[
\text{architecture observed} \Longrightarrow \text{mechanism predicted}
\]
\[
\Longrightarrow \text{player study tests the prediction}.
\]

\section{Coordination as Structural Incentive, Not Established Population-Level Outcome}

This structural account converges with an older strand of game-studies research on the unpaid labor performed by guild and alliance leadership in earlier large-scale multiplayer games \cite{koster2005}, in which leading a large player collective routinely crosses from recreational play into unpaid, intensely time-consuming work, a phenomenon related to, though not identical with, what has been termed ``playbour'' in analyses of unpaid labor in game modding \cite{kucklich2005}. The comparison should be drawn carefully rather than overstated: earlier raid-centered games were persistent, but much of their highest-coordination content could be scheduled by the participating group itself, through socially negotiated session boundaries and lockouts, whereas in globally synchronous territory games, consequential adversarial actions can instead be initiated by opponents during the defending group's absence, on a timetable the defenders do not control. It is also worth being explicit about the limits of the available evidence: the playbour literature supports an analogy between guild leadership and unpaid labor, but it does not, by itself, establish population-level empirical claims about how commonly severe or chronic consequences occur among players of the specific genre this chapter examines. This chapter's claim is structural rather than epidemiological.

\section{Remedies and Their Enforcement Difficulties}

The formal apparatus yields a design criterion in addition to a diagnosis. Let player $i$ choose, at some time $t_0$ of their own selection, an absence interval of duration $\Delta$, denoted $\mathcal{A}_i(t_0,\Delta)$. A system provides a genuine right to absence if, subject only to reasonable frequency and duration limits, $i$ can invoke $\mathcal{A}_i$ unilaterally and thereby obtain $X_i([t_0, t_0+\Delta]) \leq \varepsilon$. The word doing the essential work in this criterion is unilaterally. If invoking the interval requires permission from an alliance leader, negotiation with an adversary, or merely hoping that no one attacks during the window, the system has not supplied a right to absence; other participants are supplying protection contingently, a different and considerably weaker thing. This yields a simple diagnostic question: can the participant press leave and have the consequences stop? For a modem-era match, the answer is essentially yes. For the architecture this chapter has examined, the answer is often no, or no with qualification, particularly for participants occupying the high-$R_i$, high-$U_i$ coordinating roles the case study identified. This formalization is considerably stronger than generic ``take a break'' prompts, since a break reminder that fires while the underlying game mechanics simultaneously threaten a participant's alliance with territorial loss for having heeded it does not satisfy the criterion; the reminder is a suggestion layered on top of an architecture that still penalizes compliance, not an actual instance of $\mathcal{A}_i$.

Several distinct remedy mechanisms could contribute, each carrying its own limitation. Individual protection windows reduce a single participant's personal exposure but do not, by themselves, reduce the collective responsibility a coordinating role still bears during an active event. Alliance-level vulnerability windows address the collective case but raise the question of whose local time such a window should be anchored to when the alliance itself spans every time zone. Asynchronous attack declarations, requiring an aggressor to announce intent before a delay period elapses, could decouple event timing from any single participant's sleep schedule, at some cost to the immediacy that gives adversarial play its stakes. Hard caps on the frequency of consecutive high-stakes events limit cumulative exposure without eliminating any single event's intensity. Reversible or insured offline losses reduce $L_i(t)$ directly. Role-redundancy requirements would directly lower $U_i(t)$ for critical roles. None is a complete solution on its own. The regulatory dimension is correspondingly less settled than the loot box case, since temporal extraction has no equivalent transactional moment for a regulator to examine, and would more plausibly require product-design standards or an analogy to occupational scheduling protections rather than an existing legal category.

\section{Conclusion}

The critical and regulatory apparatus built around loot boxes has done real work, but it was built to answer a specific question, whether a randomized purchase mechanic functions as disguised gambling, and it has no purchase on a separate question: whether a game's temporal and social architecture, quite apart from anything it sells, can remove a participant's unilateral control over when disengagement is safe, and whether the resulting coordination labor becomes an uncompensated input to the game's commercial operation. The distinctive claim this chapter has defended is not that such games consume a great deal of time, since an absorbing game someone chooses to play for many hours is not thereby exploitative, but that structural unfinishability transfers control over the boundaries of play from the participant to the platform and to the other participants inhabiting it, so that leaving becomes something that can be done to a person's standing rather than something the person can simply decide to do.
\[
\boxed{\text{Exploitability can arise from the topology of permissible disengagement.}}
\]
Stated this way, the phenomenon generalizes beyond games, to on-call work schedules, notification-driven messaging platforms, always-live financial markets, and social media systems that make absence itself visible and costly to a user's standing. Money and time are extracted through different mechanisms, diagnosed by different evidence, and addressed by different remedies.

There is, finally, a structural symmetry worth noting between the argument developed here and Chapter~\ref{ch:ch04-paracosm-trap}'s companion argument about the handling of historical record. That argument holds that a system should not let a later correction retroactively rewrite an earlier commitment, that a past claim, once made, should be preserved even as further evidence constrains how it is interpreted. This chapter holds a mirror-image position about the future rather than the past: a persistent system should not let standing membership retroactively claim unlimited authority over a participant's future availability. One argument denies a system authority over what already happened; the other denies a system authority over when participation must happen next. Both reject the same underlying architectural move, a system quietly converting continuity, whether of a record or of a membership, into an unbounded claim, on the past in one case, on the future in the other, and both are instances of the qualified non-erasure principle Chapter~\ref{ch:ch10-clio-architecture} formalized as the separation of documentary history from operative authority: history that extends does not thereby license unlimited claims on what must happen next.
